% !TEX root = ../../../main.tex
% 来源：中学数学实验教材·第二册（几何）
% 章节：圆 / 圆与直线的位置关系
% 原始文件：4.tex，行 2116-2140
% 类型：tikz
% ---
  \begin{tikzpicture}[>=latex, scale=.8]
\tkzDefPoints{0/0/O, 4/0/O', 1/0/a, 2/0/b, 5/0/c}
\tkzDefMidPoint(O,O')\tkzGetPoint{O''}
\tkzDrawCircles(O,a O,b O',c)
\tkzDefPoint(75.52:1){E}
\tkzDefPoint(75.52:2){A}
\tkzDefPoint(-75.52:1){F}
\tkzDefPoint(-75.52:2){C}
\tkzDefPointsBy[translation = from O to O'](A,C){A',C'}
\tkzDefTangent[at=A](O)  \tkzGetPoint{A''}
\tkzDefTangent[at=C](O)  \tkzGetPoint{C''}
\tkzInterLL(A,A'')(O',A') \tkzGetPoint{B}
\tkzInterLL(C,C'')(O',C')\tkzGetPoint{D}
\tkzLabelPoints[above](A,B)
\tkzLabelPoints[below](C,D)
\tkzLabelPoints[left](O)
\tkzLabelPoints[above left](E)
\tkzLabelPoints[below left](F)
\tkzLabelPoints[right](O')
\tkzDrawSegments(A,O C,O B,O' D,O' O,O')
\tkzDrawSegments[dashed](E,O' F,O')
\tkzDrawCircles[dashed](O'',O)
\tkzDrawLines[thick](A,B C,D)

  \end{tikzpicture}
